\begin{claim}{C.59}{Proved}{Localized quotient-kernel maps extend over every operator translate}
Retain the plane-loop setting, $G$, $L=\overline G$,
the original commuting group $H$, an open $G$-orbit $O$,
and the full local centralizer $\mathcal C_O$ of C.55.
Fix $x\in O$, $J=L_x$, and a compact set $D\subset O$.
Let $\mathcal V_O[D]$ be the subgroup of
$\operatorname{Homeo}(O)$ consisting of maps $a$ satisfying
\[
 \begin{gathered}
 aq=qa\quad(q\in\mathcal C_O),\qquad
 \Phi_z^x a=\Phi_z^x
 \quad\bigl(z\in\operatorname{Fix}(J),\
                   Gz\text{ nonopen and nontrivial}\bigr),\\
 \operatorname{Mov}(a):=\{u:a(u)\ne u\}
                                  \subseteq\mathcal C_O D .
 \end{gathered}                                           \tag{C.59.1}
\]
Use its actual compact-open group topology, including
inverse convergence. There is a continuous injective
group homomorphism
\[
 \begin{split}
 E_D:\mathcal V_O[D]&\longrightarrow
                C_{\operatorname{Homeo}(\R^2)}(H),\\
 E_D(a)|_{hO}&=h a h^{-1}\quad(h\in H),\\
 E_D(a)|_{\R^2\setminus\bigcup_{h\in H}hO}&=1 .
 \end{split}                                             \tag{C.59.2}
\]
Restriction to $O$ is its continuous inverse onto its
image. Thus every continuous compact parameter family
in $\mathcal V_O[D]$ extends jointly with its inverse,
even when $O/\mathcal C_O$ is noncompact and infinitely
many operator translates accumulate.

The identity path component of the kernel of all the
quotient equations in (C.59.1) is trivial when the
global family has two distinct fiber partitions.
The localized perturbations in that component can
therefore contribute only in the zero-partition and
one-partition cases. The assertion does not identify
the entire quotient kernel with its identity path component.
\end{claim}

\begin{proof}
\emph{The actual group and the transported maps.}
The set $\mathcal C_OD$ is closed: the proper action
in C.55 has a Hausdorff quotient, the image of $D$
in it is compact and closed, and its full inverse
image is $\mathcal C_OD$.
The support condition is exactly pointwise fixing of
the complement. It is closed in the local homeomorphism
group, as are the centralizer and quotient conditions.
They define a closed subgroup. In particular,
$\operatorname{Mov}(a^{-1})=\operatorname{Mov}(a)$.
The inverse satisfies every quotient equation by
composing $\Phi_z^x a=\Phi_z^x$ with $a^{-1}$.

Distinct $H$-translates of $O$ are disjoint open
$G$-orbits, because each operator commutes with $G$.
If $h_1O=h_2O$, the restriction of $h_2^{-1}h_1$
lies in the actual image
$H_O\to\mathcal C_O$ in C.55.1. Equation (C.59.1)
makes its conjugation of $a$ trivial.
Thus the first rule in (C.59.2) is independent of
the chosen operator. The kernel of $H_O\to\mathcal C_O$,
which fixes $O$ pointwise, causes no additional
identification. Algebraically these rules give a
bijection of the plane, inverted by the same rules
for $a^{-1}$. They respect products and commute with
every original operator, with its original label.

\emph{The compact-set limit at every frontier point.}
Let $K$ be a compact metric parameter space and
$\lambda\mapsto a_\lambda$ a continuous map
$K\to\mathcal V_O[D]$. Write $U=\bigcup_{h\in H}hO$.
Joint continuity inside $U$ follows on each of its
open orbits. Consider
\[
 \lambda_n\to\lambda,\qquad
       u_n\in U,\qquad u_n\to z\notin U.                \tag{C.59.3}
\]
Terms fixed by $E_D(a_{\lambda_n})$ already have the
required limit. For the remaining terms choose
\[
 u_n=h_nq_n y_n,\qquad h_n\in H,\quad
             q_n\in\mathcal C_O,\quad y_n\in D .
                                                               \tag{C.59.4}
\]
After a subsequence, $y_n\to y\in D$.
The open orbit evaluation $L\to O$ supplies
$l_n\to1$ with $l_ny=y_n$.
By C.55.3 each local $q_n$ commutes with $L|_O$,
so
\[
 m_n=h_nq_n:O\longrightarrow\R^2
 \quad\hbox{is $L$-equivariant},\qquad
                     m_n(y)=l_n^{-1}u_n\to z.          \tag{C.59.5}
\]
The evaluation homeomorphism C.42.1, now based at $y$,
therefore gives
$m_n\to\Phi_z^y$ uniformly on compact subsets of $O$.
Its range is exactly $Gz$.
That orbit is nonopen: an open orbit containing $z$
would be a neighborhood disjoint from every distinct
open orbit in (C.59.3).

For a point orbit the limit map is constant.
Otherwise it is one of the actual global quotient maps,
after the exact change of source basepoint.
Namely choose $l\in L$ with $lx=y$; then
\[
                 \Phi_z^y=\Phi_{l^{-1}z}^x.             \tag{C.59.6}
\]
The point $l^{-1}z$ is fixed by $J$ and has the same
nonopen orbit type, so (C.59.1) applies.
The points $a_{\lambda_n}(y_n)$ converge to $a_\lambda(y)$
and lie in the compact image of $K\times D$.
Consequently (C.59.4)--(C.59.6) imply
\[
 \begin{split}
 E_D(a_{\lambda_n})(u_n)
     &=h_nq_n a_{\lambda_n}(y_n)\\
     &\longrightarrow\Phi_z^y(a_\lambda(y))
       =\Phi_z^y(y)=z .
 \end{split}                                             \tag{C.59.7}
\]
Every subsequence admits the same limit argument,
so this proves convergence for the entire sequence.
Sequences outside $U$ are fixed.
This proves joint continuity on the whole plane.

The inverse family is continuous locally and has
the identical moving-point saturation by (C.59.1).
The same proof gives joint continuity of its transported
family. Thus the algebraic inverse above is a continuous
inverse, and every transported map is a homeomorphism.
Compact parameter families give continuous maps into
the ambient compact-open group.
For continuity of $E_D$ itself, apply this conclusion
to a convergent sequence and its limit in
$\mathcal V_O[D]$, which is a compact metrizable
parameter set. Both homeomorphism groups are metrizable,
so sequential continuity proves continuity.
Restriction to $O$ is continuous and recovers $a$,
proving the claimed topological embedding.
No compactness of the entire quotient was used;
only the one compact set $D$ in (C.59.4) was used.

\emph{The exact two-partition restriction.}
Let $a_s$, $0\leq s\leq1$, be a path from the identity
fixing all global quotient maps, and choose two maps
with distinct partitions.
Retain their original source and target covers
$p:\widetilde O\to O$ and $q_i:\R\to B_i$.
The corresponding real lifts $F_i$ satisfy
$q_iF_i=\Phi_i p$.
C.45 proves that
$(F_1,F_2)$ is injective on $\widetilde O$.
Lift the isotopy $a_s p$ to the based source isotopy
$\widetilde a_s$, with $\widetilde a_0=1$.
For each $i$ the quotient equation gives
\[
                q_iF_i\widetilde a_s=q_iF_i .          \tag{C.59.8}
\]
For a line target this is equality of the two real
maps. For a circle their difference is an integer
deck translation, continuous in $(s,u)$ and zero
at $s=0$. It is therefore zero throughout.
Thus both real coordinates of $\widetilde a_s(u)$
equal those of $u$. Pair injectivity makes
$\widetilde a_s=1$, and descent gives $a_s=1$.
This retains the discrete deck ambiguity for maps
outside the identity path component.

For representatives outside the invariant closed set
$X^\flat$, (C.59.2) fixes $X^\flat$ pointwise and gives
relative perturbations there. It does not fragment
the original loop: on a transitive open orbit the
union of that loop's moving sets is the whole orbit,
so imposing one compact quotient support on it would
force a compact quotient. The new assertion supplies
the exact extension for localized factors; existence
and convergence of a factorization sufficient for
C.12 are still unproved.
\end{proof}
